Throughout $n(x) = 1/(e^{x} + 1)$ and $s(x) = \ln(1 + e^{-x}) + xn(x)$ is the entropy of one mode (Theorem \claimref{synthesis.boundary-mode-thermal-entropy}), even in $x$ with $s(0) = \ln 2$. The spectral form of the continuum is the stored one: $h$ acts as $-\kappa w$ on $u^{-1/2 + I_4w}$ and the symbol of the I-state is $(1 + e^{\beta h})^{-1}$ (Theorems \claimref{synthesis.positive-null-energy-gives-r-premises}(ii) and \claimref{hawking.i-state-continuum}); the Mellin transform identifies the one-particle space with $L^{2}(\R, d\varpi)$, $\varpi$ the eigenvalue of $h$.

\emph{(i)} The intervals $\tfrac{\kappa}{r}[k, k + 1)$, $k \in \Z$, partition $\R$, so $\int_{\R}|\varphi|^{2}d\varpi = \sum_{k}\int_{0}^{1}|\varphi(\tfrac{\kappa}{r}(k + \delta))|^{2}\tfrac{\kappa}{r}\,d\delta$, which is $\|U\varphi\|^{2}$ by Tonelli's theorem; $U$ is onto, with inverse $\varphi(\tfrac{\kappa}{r}(k + \delta)) = \sqrt{r/\kappa}\,\Phi(\delta)_{k}$. A multiplication operator $F(\varpi)$ is carried to the family $\operatorname{diag}(F(\tfrac{\kappa}{r}(k + \delta)))$: for $F(\varpi) = \varpi$ this is $h_{\delta}$; for $F = (1 + e^{\beta\varpi})^{-1}$ it is $\operatorname{diag}(1/(e^{x_{k}} + 1))$, the symbol of the Gibbs product state of $\mathcal{A}(\delta, r)$ (which is the quasi-free KMS state at $\beta$ of the flow generated by $h_{\delta}$, by Theorem \claimref{hawking.i-state-continuum}(1) in its point-spectrum case); for $F = e^{I_4L\varpi}$, $L\tfrac{\kappa}{r}(k + \delta) = 2\pi(k + \delta)$ gives $e^{2\pi I_4\delta}$, independent of $k$. The phase is $-1$ exactly at $\delta = \tfrac12$ and $+1$ exactly at $\delta = 0$. If $e^{I_4L\varpi}\varphi = \lambda\varphi$ almost everywhere, $\varphi$ vanishes off the countable set where $e^{I_4L\varpi} = \lambda$, so $\varphi = 0$: no eigenvector. A vector in a single fibre would be an eigenvector with $\lambda = e^{2\pi I_4\delta}$. A subalgebra $\CAR(\mathfrak{k})$, $\mathfrak{k}$ a closed subspace invariant under $h$, on which the restricted I-state were the Gibbs state of $\mathcal{A}(\delta, r)$ with the Killing flow as its flow would need $h|_{\mathfrak{k}}$ of pure point spectrum, and $h$ has no eigenvectors (stored, clause (iv) of the cited theorem). This is the exact form of that clause's statement that the ladder values are recovered only for modes of vanishing Mellin width. Without invariance the conclusion fails: let $I_{k}$ be disjoint open intervals with $\tfrac{\kappa}{r}(k + \delta) \in I_{k}$. The symbol $m(\varpi) = (1 + e^{\beta\varpi})^{-1}$ is continuous and strictly monotone, and $\int|\phi|^{2}m\,d\varpi$ over unit vectors supported in $I_{k}$ takes every value strictly between the infimum and supremum of $m$ on $I_{k}$ (move the weight of $\phi$ continuously across the interval), in particular $m(\tfrac{\kappa}{r}(k + \delta)) = 1/(e^{x_{k}} + 1)$. For such $\phi_{k}$, orthonormal by disjoint supports, the compression of the symbol to their closed span is $\operatorname{diag}(1/(e^{x_{k}} + 1))$ (off-diagonal entries vanish by disjoint supports), so the restricted quasi-free state is the Gibbs product state of $\mathcal{A}(\delta, r)$. The span is not invariant under $h$, and the construction works for every $\delta$ and $r$.

\emph{(ii)} For $\varphi, \chi$ absolutely continuous on $[0, L]$ with square-integrable derivatives, integration by parts gives $\ip{\chi}{I_4\varphi'} - \ip{I_4\chi'}{\varphi} = I_4\bigl(\overline{\chi(L)}\varphi(L) - \overline{\chi(0)}\varphi(0)\bigr)$, the bar denoting the conjugation of $\mathrm{span}\{1, I_4\}$. The operator on the domain $\varphi(L) = e^{-2\pi I_4\delta}\varphi(0)$ is symmetric (the boundary form vanishes) and has the complete orthonormal eigenbasis $L^{-1/2}e^{-I_4\varpi_{k}\tau}$, $\varpi_{k} = \tfrac{2\pi}{L}(k + \delta) = \tfrac{\kappa}{r}(k + \delta)$ (the twisted Fourier basis; $I_4\,d/d\tau$ has eigenvalue $\varpi_{k}$ on it and $\varphi(L) = e^{-2\pi I_4\delta}\varphi(0)$), so it is self-adjoint with that spectrum. Conversely the minimal operator, with $\varphi(0) = \varphi(L) = 0$, has deficiency indices $(1, 1)$ (the solutions $e^{\pm\tau}$ of $I_4\varphi' = \pm I_4\varphi$, one each), so by von Neumann's theorem its self-adjoint extensions form one circle, and each is a restriction of the maximal operator on which the boundary form vanishes, which forces $|\varphi(L)| = |\varphi(0)|$ with a fixed phase: the twisted conditions exhaust them. The Gibbs state of $\mathcal{A}(\delta, r)$ is quasi-free and KMS at $\beta$ for the flow of $h_{\delta}$ for every $\delta$. The relabelling $k \mapsto -1 - k$ (or $k \mapsto -k$) preserves the weight set up to sign exactly when $2\delta \in \Z$; the two classes and their flows at the recurrence time are stored (Theorem \claimref{synthesis.q-splits-maximum-entropy-affine-weights}(iv)). The cited theorem's clause (iii) states that the weight clause of (Q) is not given by (N$_{+}$).

\emph{(iii)} Write $k = Mm + j$ with $m \in \Z$, $0 \le j < M$: $\tfrac{\kappa}{Mr}(k + \delta) = \tfrac{\kappa}{r}\bigl(m + \tfrac{j + \delta}{M}\bigr)$, a bijection between the modes of $\mathcal{A}(\delta, Mr)$ and the disjoint union of the modes of the $M$ arrays. A Gibbs product state is the product over modes of one-mode states that depend only on $x$, so the state of $\mathcal{A}(\delta, Mr)$ is the product of the $M$ Gibbs states (for even states the tensor product is the product in the algebra), and entropy, energy above the sea and variance, each a sum over modes, add. For $\delta = \tfrac12$, $\delta_{j} = \tfrac{2j + 1}{2M}$ and $e^{2\pi I_4\delta_{j}M} = e^{\pi I_4(2j + 1)} = -1$. With $\sum_{j}(2j + 1) = M^{2}$ and $\sum_{j}(2j + 1)^{2} = M(4M^{2} - 1)/3$, $\sum_{j}B_{2}(\delta_{j}) = \tfrac{4M^{2} - 1}{12M} - \tfrac{M}{2} + \tfrac{M}{6} = -\tfrac{1}{12M} = B_{2}(\tfrac12)/M$, so $\sum_{j}\tfrac{\pi}{r}B_{2}(\delta_{j}) = \tfrac{\pi}{Mr}B_{2}(\tfrac12)$, the sea of the long window. For a general twist, with $\sum_{j}j = M(M - 1)/2$ and $\sum_{j}j^{2} = M(M - 1)(2M - 1)/6$, the same computation gives $\sum_{j}B_{2}\bigl(\tfrac{j + \delta}{M}\bigr) = (\delta^{2} - \delta + \tfrac16)/M$. The case $M = 2$, $r = \tfrac12$: twists $\tfrac14, \tfrac34$, exchanged by $k \mapsto -1 - k$, so the two factors have equal entropy, energy and sea, half of $\tfrac12\ln 2$ (Theorem \claimref{synthesis.half-integer-weight-array-carries-half-ln2}(i)), of $\pi/12$ and of $-\pi/12$; directly, $\tfrac{\pi}{1/2}B_{2}(\tfrac14) = 2\pi\cdot(-\tfrac{1}{48}) = -\pi/24$.

\emph{(iv)} By Tonelli's theorem and the substitution $t = k + \delta$, $\int_{0}^{1}\sum_{k}g(2\pi(k + \delta)/r)\,d\delta = \int_{\R}g(2\pi t/r)\,dt = \tfrac{r}{2\pi}\int_{\R}g$. Apply it to $g = s$, $g(x) = |x|n(|x|)$ and $g(x) = x^{2}n(x)(1 - n(x))$, all nonnegative, even and integrable (a mode of weight zero occurs only at $\delta = 0$, a null set). \emph{The values.} By clauses (a), (b) of Fact \claimref{synthesis.imported-theta-function-identities}, $\vartheta_{3}/\eta$ is invariant under $\tau \mapsto -1/\tau$, which at $\tau = I_4/R$ reads $\tfrac{\pi}{12R} + 2\sum_{n \ge 1}\ln(1 + e^{-2\pi a_{n}/R}) = \tfrac{\pi R}{12} + 2\sum_{n \ge 1}\ln(1 + e^{-2\pi Ra_{n}})$, $a_{n} = n - \tfrac12$. Applying $1 + R\,d/dR$, $R\,d/dR$ and $(R\,d/dR)^{2} + R\,d/dR$ termwise gives, with $y_{n} = 2\pi Ra_{n}$, $S_{1/2}(R) = \tfrac{\pi R}{6} + 2\sum_{n}[\ln(1 + e^{-y_{n}}) - y_{n}n(y_{n})]$, $e_{1/2}(R) = \tfrac{\pi R}{12} + \tfrac{\pi}{12R} - 2\sum_{n}y_{n}n(y_{n})$ and $v_{1/2}(R) = \tfrac{\pi R}{6} + 2\sum_{n}y_{n}n(y_{n})[y_{n}(1 - n(y_{n})) - 2]$; each sum is bounded by $4\sum_{n}(1 + y_{n})^{2}e^{-y_{n}}$, which tends to $0$ as $R \to \infty$. (This is the dual form of the base comparison of this topic, repeated here.) On the other hand $S_{1/2}(R)/R = \tfrac{1}{2\pi}\sum_{k}\epsilon\,s(\epsilon(k + \tfrac12))$ with $\epsilon = 2\pi/R$ is a Riemann sum, and for an integrable function of bounded variation the Riemann sum differs from the integral by at most $\epsilon$ times the total variation; $s$, $|x|n(|x|)$ and $x^{2}n(1 - n)$ are integrable and of bounded variation. Hence $\tfrac{1}{2\pi}\int_{\R}s = \tfrac{\pi}{6}$, $\tfrac{1}{2\pi}\int_{\R}|x|n(|x|) = \tfrac{\pi}{12}$, $\tfrac{1}{2\pi}\int_{\R}x^{2}n(1 - n) = \tfrac{\pi}{6}$, and the three averages follow. $\int_{0}^{1}B_{2} = \tfrac13 - \tfrac12 + \tfrac16 = 0$. The continuum's energy over the window in the boost ordering is $\pi r/12$ by clause (i) of Theorem \claimref{synthesis.null-ordering-of-counted-units-is-zeta-symmetric} (density $\kappa^{2}/(48\pi)$ times $L\beta$). The $M$-point statement is (iii) divided by $M$, and the values are $\Delta(M)/M$ with $\Delta(R) = S_{1/2}(R) - \pi R/6$, bounded by the check.

\emph{(v)} The relabelling $k \mapsto -1 - k$ maps the weights of twist $\delta$ to minus those of twist $1 - \delta$, and $s$ is even. The values are certified by rational interval arithmetic in the check: $S_{0}$ is the stored value of Theorem \claimref{synthesis.half-integer-weight-array-carries-half-ln2}(ii) (the $\ln 2 + \varepsilon$ of Theorem \claimref{synthesis.residue-from-mode-entropy}(vi)), re-bounded. $S_{\delta}(1)$ is continuous in $\delta$ on $(0, 1)$ (a locally uniformly convergent sum of continuous functions), and $S_{6/25} > \pi/6 > S_{1/4}$, so it takes the value $\pi/6$ in between. A twist is particle-hole symmetric only if $2\delta \in \Z$.

\emph{Remark.} The antiperiodic window mode $k$ at $r = 1$ is $\phi_{k}(\tau) = L^{-1/2}e^{-I_4\varpi_{k}\tau}$ on $[0, L]$ and zero elsewhere; its spectral density in $h$ is $\tfrac{1}{2\pi L}\,4\sin^{2}((\varpi - \varpi_{k})L/2)/(\varpi - \varpi_{k})^{2}$, and in $x = \beta\varpi$ with $L = \beta$, $\sin^{2}((x - x_{k})/2) = \cos^{2}(x/2)$ since $x_{k} = (2k + 1)\pi$. Its occupation in the I-state is the integral of this density against $n(x)$ (clause (iv) of the cited theorem), the stated $n_{k}$; the numerical values are from a floating-point quadrature and are not evidence in the check. For large $k$ the part $x < 0$, where $n \to 1$, gives $\tfrac{2}{\pi}\int_{-\infty}^{0}\cos^{2}(x/2)(x - x_{k})^{-2}dx \approx \tfrac{1}{\pi x_{k}}$. If the restricted state were given by a density matrix in the Fock space over the array's sea, then by the Powers and Stormer criterion (from memory) $\sum_{k \ge 0}n_{k}$ would be finite. This is argued, not proved here.

\emph{Status.} (i), (iii) and the identity of (iv) are proved and elementary; the values in (iv) are proved given the imported Fact; (ii) is classical with the boundary form computed; (v) is certified by exact bounds, with every value between the two end values taken by continuity and extremality not proved; the Remark is argued. That the clauses do not depend on Premise (Q), (N$_{+}$) or the window reading is argued at clause level; the declared cone contains them through the cited theorems.
